\documentclass[12pt,a4paper]{article}
\usepackage[utf8]{inputenc}
\usepackage[vietnamese]{babel}
\usepackage[left=2cm,right=2cm,top=2cm,bottom=2cm]{geometry}
\usepackage{mathptmx}
\usepackage{enumitem}

\pagestyle{plain}
\linespread{1.15}

\begin{document}

% --- HEADER QUỐC HIỆU ---
\noindent
\begin{minipage}[t]{0.45\textwidth}
    \begin{center}
        \small
        CÔNG AN THÀNH PHỐ HÀ NỘI\\
        \textbf{\underline{PHÒNG CẢNH SÁT HÌNH SỰ}}\\[0.1cm]
        Số: 11/BB-TXCAM\\
        \textit{V/v khám xét nơi ở và lời khai tự thú của Tùng}
    \end{center}
\end{minipage}
\hfill
\begin{minipage}[t]{0.52\textwidth}
    \begin{center}
        \small
        \textbf{CỘNG HÒA XÃ HỘI CHỦ NGHĨA VIỆT NAM}\\
        \textbf{\underline{Độc lập - Tự do - Hạnh phúc}}\\[0.1cm]
        \textit{Hà Nội, ngày 25 tháng 07 năm 2026}
    \end{center}
\end{minipage}

\vspace{0.6cm}

\begin{center}
    \textbf{\large BIÊN BẢN KHÁM XÉT NƠI Ở VÀ LỜI KHAI TỰ THÚ}\\
    \textit{(Làm rõ hành vi xô xát của nghi phạm Nguyễn Thanh Tùng)}
\end{center}

\vspace{0.3cm}

\section*{I. KẾT QUẢ KHÁM XÉT PHÒNG RIÊNG CỦA TÙNG (KÝ HIỆU f3-2)}
\noindent Căn cứ Lệnh khám xét khẩn cấp số 18/LKX-CSHS;\\
Tại phòng trọ của Tùng, cơ quan điều tra thu giữ 01 hộp các-tông đặt trên nóc tủ:
\begin{itemize}[leftmargin=1.5cm]
    \item 01 Khung tranh kính cũ vẽ hai anh em nắm tay nhau. Tháo mặt sau khung tranh phát hiện dòng chữ viết tay phai màu mực tím: \textit{"Thương em, anh xin lỗi vì đã không tìm thấy được em. An nghỉ em nhé. Hè 1998 — Anh Tùng"}.
    \item Các tài liệu cá nhân: Sổ tay ghi chép công trường, biên lai điện nước và đồng hồ cũ.
\end{itemize}

\section*{II. TRÍCH XUẤT CAMERA VÀ MẢNH BÁO BỊ XÉ TẠI HIỆN TRƯỜNG}
\begin{itemize}[leftmargin=1.5cm]
    \item \textbf{Camera cây xăng Petrolimex số 12 (cách 300m):} Ghi nhận Tùng vào ngõ lúc 19:45 và hốt hoảng tháo chạy lúc **20:15** bằng xe máy.
    \item \textbf{Ghép hoàn chỉnh mảnh báo cũ năm 1998:} Tiêu đề \textit{"Sự cố thương tâm: Bé trai tử vong do ngạt khí khi bị kẹt trong tủ"} — làm rõ bé Gia Huy (bị câm, bệnh tim) chết ngạt do chốt gỗ bên ngoài sập xuống.
\end{itemize}

\section*{III. LỜI KHAI TỰ THÚ TOÀN BỘ CỦA NGUYỄN THANH TÙNG}
\noindent Trước các chứng cứ không thể chối cãi (Ghim cài áo, Bút tích khung tranh, Mảnh báo rách), Nguyễn Thanh Tùng đã khai nhận toàn bộ:
\begin{itemize}[leftmargin=1.5cm]
    \item \textbf{Nguồn cơn:} Hai hôm trước, trong bữa nhậu say, Khang cười cợt khoe chiến tích thuở nhỏ: \textit{"Đỉnh nhất là cái vụ gài chốt nhốt thằng nhóc vào tủ, tìm hoài không ra!"}. Tùng bàng hoàng nhận ra Khang chính là kẻ đã cố tình nhốt chết em trai mình.
    \item \textbf{Xô xát lúc 20:00:} Tùng cầm bài báo sang nhà Khang đối chất. Khang thờ ơ nhạo báng rồi giật xé nát bài báo. Quá phẫn nộ, Tùng giằng co làm vỡ bộ bình trà và xô ngã Khang đập đầu vào tủ ngất xỉu.
    \item \textbf{Tháo chạy lúc 20:15:} Tùng kiểm tra thấy Khang **vẫn còn thở**, nhưng vì quá hoảng sợ nên đã bỏ chạy khỏi hiện trường lúc 20:15 để bắt xe khách về Hải Phòng.
\end{itemize}

\vspace{0.8cm}

\noindent
\begin{minipage}[t]{0.45\textwidth}
    \small
    \textbf{\textit{Nơi nhận:}}\\
    - Trưởng Ban Chuyên án;\\
    - Lưu: VT, CSHS.
\end{minipage}
\hfill
\begin{minipage}[t]{0.5\textwidth}
    \begin{center}
        \small
        \textbf{ĐIỀU TRA VIÊN CHỦ TRÌ}\\
        \textit{(Ký, ghi rõ họ tên)}\\[1.5cm]
        \textbf{Đại úy Lê Minh}
    \end{center}
\end{minipage}

\end{document}
